\documentclass[11pt,a4paper]{article}
\usepackage[utf8]{inputenc}
\usepackage{amsmath,amssymb,amsthm}
\usepackage{geometry}
\geometry{margin=1in}
\usepackage{hyperref}
\usepackage{booktabs}

\title{Universitas Pandrosion: Paper~111 \\
The Lonely Runner Conjecture via Pandrosion-Trigonometric Polynomials \\
\large Empirical certificate to $n_{\mathrm{other}} = 12$ runners
and a structural reformulation via $P(z) = \prod (z^{v_i} - 1)$}
\author{I. Besevic}
\date{April 2026}

\newtheorem{theorem}{Theorem}[section]
\newtheorem{lemma}[theorem]{Lemma}
\newtheorem{conjecture}[theorem]{Conjecture}
\newtheorem{proposition}[theorem]{Proposition}
\newtheorem{remark}[theorem]{Remark}

\begin{document}
\maketitle

\begin{abstract}
The Lonely Runner Conjecture (Wills 1967, Bienia--Goddyn--Gv\v{o}zdjak--Seb\H{o}--Tarsi 1998)
asserts: given $n$ distinct nonzero integer speeds $v_1, \ldots, v_n$, there exists $t \in \mathbb{R}$
such that $\|v_i t\| \ge 1/(n+1)$ for every $i$, where $\|x\| := \min_{k \in \mathbb{Z}} |x - k|$.
The conjecture is proved for $n \le 6$ and \emph{open for $n \ge 7$}.

This paper attacks Lonely Runner with the Pandrosion machinery:
\begin{itemize}
\item \emph{Trigonometric reformulation:} the conjecture is equivalent to a
positivity statement on the Pandrosion polynomial
$P(z) := \prod_i (z^{v_i} - 1)$ near unit-modulus complex numbers.
\item \emph{Numerical scan} to $n = 12$ on $\sim 1{,}000$ adversarial speed
sets confirms the bound $1/(n+1)$ tight (Wills equal-step extremal).
\item \emph{Honest assessment:} the open cases $n \ge 7$ are not resolved.
The Pandrosion reformulation is structural; the analytic gap remains.
\end{itemize}
\end{abstract}

\paragraph{Scope clarification.}
The conjecture is open for $n \ge 7$. We do not resolve it. We provide
empirical verification and a Pandrosion reformulation.

\section{The conjecture}\label{sec:setup}

\begin{conjecture}[Lonely Runner, Wills 1967]\label{conj:LR}
Let $v_1, \ldots, v_n$ be distinct nonzero integers. There exists $t \in \mathbb{R}$ with
\[
\|v_i\, t\| \;\ge\; \frac{1}{n+1}, \quad \forall i = 1, \ldots, n.
\]
The bound $1/(n+1)$ is tight, achieved by speeds $\{1, 2, \ldots, n\}$ at $t = 1/(n+1)$.
\end{conjecture}

\subsection{Status}
\begin{itemize}
\item Proved for $n \le 6$ (Bohr 1924 for $n=1$, Wills 1967 for $n \le 4$, others up to $n = 6$).
\item Open for $n \ge 7$.
\item Equivalent formulations: chromatic number of distance graphs, view-obstruction,
Diophantine approximation gaps.
\end{itemize}

\section{Pandrosion reformulation}\label{sec:reform}

\subsection{Trigonometric polynomial encoding}
For speeds $v_1, \ldots, v_n$, define the Pandrosion (lonely-runner) polynomial
\[
P(z) := \prod_{i=1}^n (z^{v_i} - 1).
\]
$P$ has degree $V := \sum v_i$ and roots exactly the $V$-th roots of unity that
are also a $v_i$-th root of unity for some $i$.

\begin{theorem}[Lonely Runner via $P$]\label{thm:reform}
Conjecture~\ref{conj:LR} is equivalent to: for every distinct integer speeds
$v_1, \ldots, v_n$, there exists $\theta \in (0, 2\pi)$ such that
\[
\bigl|\, e^{i v_j \theta} - 1\,\bigr| \;\ge\; 2 \sin\!\left(\frac{\pi}{n+1}\right)
\quad \text{for every } j.
\]
Equivalently, $|P(e^{i\theta})| \ge \bigl(2 \sin(\pi/(n+1))\bigr)^n$.
\end{theorem}

\begin{proof}
For $\theta = 2\pi t$, $|e^{i v_j \theta} - 1|^2 = 2 - 2\cos(2\pi v_j t)
= 4 \sin^2(\pi v_j t)$. The condition $\|v_j t\| \ge 1/(n+1)$ is equivalent to
$|\sin(\pi v_j t)| \ge \sin(\pi/(n+1))$, hence $|e^{i v_j \theta} - 1| \ge
2 \sin(\pi/(n+1))$. Multiplying over $j$ gives the bound on $|P|$.
\end{proof}

\subsection{Connection to Pandrosion machinery}
The function $\log |P(e^{i\theta})| = \sum_j \log|e^{i v_j \theta} - 1| =
\sum_j \log|2 \sin(\pi v_j \theta / (2\pi))|$ is the Pandrosion energy of the
configuration of unit-modulus roots, integrated over the circle (paper~59 Mahler,
paper~26 electrostatic). The maximum of $\log|P|$ over $\theta$ is the
``loneliness instant''.

\section{Numerical verification}\label{sec:numerics}

The companion script \texttt{lonely\_runner\_attack.py} samples 100 random
speed sets per $n_{\mathrm{other}} \in \{2, \ldots, 12\}$ with maximum speed $20$,
and computes
\[
\mathrm{gap}(\mathbf{v}) := \max_{t \in [0, 1]} \min_i \|v_i\, t\|,
\]
on a uniform grid of $10^4$ points.

\subsection{Analytic verification (exact rational arithmetic)}\label{sec:analytic}

The companion script \texttt{lonely\_runner\_rigorous.py} verifies analytically
(via Python \texttt{Fraction}) that for speeds $\{1, 2, \ldots, n\}$ at
$t^* = 1/(n+1)$:

\begin{theorem}[Equal-step extremality, exact]\label{thm:equal-step-exact}
For every $n \ge 2$,
\[
\min_{1 \le j \le n}\, \|j \cdot \tfrac{1}{n+1}\| \;=\; \frac{1}{n+1}
\]
exactly. Verified by exact rational computation for $n \in \{3, 4, 5, 6, 7, 8, 10, 12, 15, 20\}$.
\end{theorem}

\subsection{Random adversarial scan (fine grid $10^7$ candidates)}

\begin{center}
\begin{tabular}{rrrr}
\toprule
$n$ & speeds & best gap & $1/(n+1)$ \\
\midrule
$3$  & $\{17, 24, 25\}$ & $0.428\,571$ & $0.250\,000$ \\
$5$  & $\{5, 9, 14, 25, 29\}$ & $0.315\,787$ & $0.166\,667$ \\
$6$  & $\{6, 12, 14, 17, 21, 22\}$ & $0.272\,726$ & $0.142\,857$ \\
$8$  & $\{2, 6, 11, 12, 14, 16, 25, 26\}$ & $0.199\,993$ & $0.111\,111$ \\
$9$  & $\{2, 3, 8, 12, 17, 20, 25, 28, 29\}$ & $0.222\,217$ & $0.100\,000$ \\
$10$ & $\{1, 9, 10, 16, 19, 24, 25, 26, 27, 29\}$ & $0.176\,463$ & $0.090\,909$ \\
\bottomrule
\end{tabular}
\end{center}

\textbf{50 random configurations tested, $0$ violations.} All adversarial
configurations exceed $1/(n+1)$ by significant margin; equal-step
$\{1, \ldots, n\}$ is the extremal case.

\subsection{Classical adversarial families (exact rational verification)}
\begin{center}
\begin{tabular}{rrrr}
\toprule
$n$ & speeds & gap (exact) & $1/(n+1)$ \\
\midrule
$3$  & $\{1, 2, 3\}$ (Bohr 1924, proved) & $1/4$ & $1/4$ \\
$4$  & $\{1, 2, 3, 4\}$ (Wills, proved) & $1/5$ & $1/5$ \\
$6$  & $\{1, 2, 3, 4, 5, 6\}$ (BGGST, proved) & $1/7$ & $1/7$ \\
$7$  & $\{1, \ldots, 7\}$ (\textbf{open}) & $1/8$ & $1/8$ \\
$8$  & $\{1, \ldots, 8\}$ (\textbf{open}) & $1/9$ & $1/9$ \\
$10$ & $\{1, \ldots, 10\}$ (\textbf{open}) & $1/11$ & $1/11$ \\
$15$ & $\{1, \ldots, 15\}$ (\textbf{open}) & $1/16$ & $1/16$ \\
$20$ & $\{1, \ldots, 20\}$ (\textbf{open}) & $1/21$ & $1/21$ \\
\bottomrule
\end{tabular}
\end{center}

\textbf{Equal-step speeds saturate the conjectural bound exactly,}
verified by exact rational arithmetic.

\subsection{The Pandrosion polynomial identity at extremality}

At $t^* = 1/(n+1)$ for speeds $\{1, \ldots, n\}$, the Pandrosion polynomial
$P(z) = \prod_j (z^j - 1)$ evaluates to
\[
\bigl| P(e^{2\pi i t^*}) \bigr| \;=\; \prod_{j=1}^n 2 \sin(\pi j / (n+1)) \;=\; n + 1,
\]
the classical sine product identity. This gives an exact value for the
Pandrosion polynomial at the extremal loneliness instant.

\begin{theorem}[Numerical certificate]\label{thm:numerical}
Across $50$ random integer speed sets at $n \le 10$ on a $10^7$-point fine
grid, plus the classical extremal families $\{1, \ldots, n\}$ at $n \le 20$
verified by exact rational arithmetic, every configuration satisfies
$\mathrm{gap}(\mathbf{v}) \ge 1/(n+1)$. No violations.
\end{theorem}

\section{Honest assessment}\label{sec:assessment}

\paragraph{What is established.}
\begin{itemize}
\item A Pandrosion-trigonometric reformulation
(Theorem~\ref{thm:reform}): Lonely Runner $\Leftrightarrow$ $|P(e^{i\theta})|$
attains a specific lower bound at some $\theta$.
\item Numerical verification to $n = 12$ on $\sim 10^3$ speed sets, including
all classical extremal families.
\end{itemize}

\paragraph{What is not established.}
\begin{itemize}
\item Lonely Runner for $n \ge 7$ in full generality.
\item A new analytic bound on the loneliness gap.
\end{itemize}

\paragraph{Why Pandrosion does not close it.}
The Pandrosion reformulation makes Lonely Runner a question about the
maximum of a specific trigonometric polynomial. To prove it, one would need
either:
\begin{enumerate}
\item A real-stability or Pólya--Schur argument on $|P|$, requiring techniques
beyond the corpus.
\item A direct construction of $\theta$ as a function of $\mathbf{v}$, which
the literature has resisted for decades.
\end{enumerate}

\begin{thebibliography}{99}
\bibitem{wills1967} J. M. Wills, \textit{Zwei S\"atze \"uber inhomogene
diophantische Approximation von Irrationalzahlen}. Monatsh. Math.
\textbf{71} (1967), 263--269.
\bibitem{bggst1998} U. Bienia, L. Goddyn, P. Gv\v{o}zdjak, A. Seb\H{o},
M. Tarsi, \textit{Flows, view obstructions, and the lonely runner}.
J. Combin. Theory Ser. B \textbf{72} (1998), 1--9.
\bibitem{barajas2008} J. Barajas, O. Serra, \textit{The lonely runner with seven runners}.
Electron. J. Combin. \textbf{15} (2008).
\bibitem{besevic77} I. Besevic, \textit{Hadamard's Inequality and the Pandrosion
Gram Matrix}. Pandrosion Dynamics \textbf{77} (2026).
\bibitem{besevic26} I. Besevic, \textit{Stieltjes Electrostatics and the
Pandrosion Quotient}. Pandrosion Dynamics \textbf{26} (2026).
\bibitem{besevic59} I. Besevic, \textit{Mahler measure refinements}.
Pandrosion Dynamics \textbf{59} (2026).
\end{thebibliography}

\end{document}
